\chapter*{Introduction}
\addcontentsline{toc}{chapter}{Introduction}

A PhD student sits in his advisor’s office, scrambling through piles of computer outputs in an attempt to finish a task his advisor assigned. A visiting professor enters, and their ensuing conversation leads to a question so intriguing that it consumes the student’s thoughts overnight. The next morning, he rushes to the lab and begins coding. This is the origin story of the GARCH model, as narrated by \citeA{Bollerslev2023}, who describes the creation of the model as a personal odyssey. At that moment, Bollerslev had little idea of the profound impact his model would have, revitalizing and significantly expanding the literature on financial return modeling. This literature quickly grew, matured, and branched out, generating numerous alternative approaches for nearly every aspect of return modeling. This thesis is, in essence, a consequence of the question posed in 1984 and the extensive literature it subsequently inspired.

The goal of this thesis is to provide a comprehensive overview of financial return modeling. While the thesis acknowledges and discusses alternative methodologies at various modeling stages, it follows a specific path designed to provide the reader with a clear, coherent, and complete narrative. Ultimately, it aims to serve as a concise yet thorough "recipe" for modeling financial returns at the undergraduate level.

Special acknowledgement is due to \citeA{Christoffersen2012}, whose textbook \emph{Elements of Financial Risk Management} served as a guiding reference throughout this work. The aim of presenting the material as a clear, coherent, and reasonably complete narrative was directly inspired by the pedagogical clarity of that text. In the author’s view, Christoffersen’s exposition fills an important gap between the mathematical and statistical foundations typically acquired at the undergraduate level and the more demanding world of quantitative finance. It was through this text that the author first came to appreciate the depth and appeal of financial mathematics.

\hyperref[chap:1]{Chapter 1} outlines the fundamental task of return modeling, highlighting common empirical characteristics that models must capture and replicate. Chapters 2 through 4 then guide the reader through the univariate modeling process. \hyperref[chap:2]{Chapter 2} addresses volatility modeling, \hyperref[chap:3]{Chapter 3} focuses on modeling the distribution of shocks, and \hyperref[chap:4]{Chapter 4} presents methods for rigorously evaluating and testing these models. \hyperref[chap:5]{Chapter 5} extends the framework to a multivariate setting, building coherently on the foundations established in the previous chapters.

All \Rlogo\ code used to generate the analyses and figures is available in the accompanying GitHub repository: \href{https://github.com/ferpmillan/thesis}{github.com/ferpmillan/thesis}.